\slajdid{09_2-s051}
\begin{frame}[label=09_2-s051]
\gusttekst\setlength{\abovedisplayskip}{3pt}\setlength{\belowdisplayskip}{3pt}\setlength{\abovedisplayshortskip}{2pt}\setlength{\belowdisplayshortskip}{2pt}\begin{columns}[c,onlytextwidth]\column{.32\textwidth}\centering\input{slike/tikz/s051_pozitivna_povratna_sprega.tex}\column{.66\textwidth}
Ono što se može uočiti sa slike je da kondenzatori, parazitne kapacitivnosti, $C_p$ i $C_N$ u odnosu na izlazni signal i baze tranzistora $T_N$ i $T_P$ imaju diferencijatorski efekat. Što je brzina porasta signala na izlazu veća to će i ovaj efekat biti izraženiji. Može se desiti da impuls na bazi tranzistora $T_N$ bude dovoljno veliki da se uključi tranzistor. Tada započinje pozitivna reakcija i nastaje lečap. Tranzistor $T_N$ kada provede izazvaće pad napona na otporniku $R_{W1}$ što dovodi do uključenja tranzistora $T_P$. Uključenje tranzistor $T_P$ izaziva pojavu napona na otporniku $R_{S1}$ što dalje obezbeđuje bolje provođenje tranzistora $T_N$, veći pad napona na otporniku $R_{W1}$, bolje provođenje tranzistora $T_P$, veći pad napona na otporniku $R_{S1}$, bolje provođenje tranzistora $T_N$… Pozitivna reakcija koja će se završiti tako što i tranzistori $T_N$ i $T_P$ vode bez ikakvih spoljnih dodatnih uslova, i ostvaruju kratak spoj između napajanja i mase. Do iste pozitivne reakcije vodi i situacija ako se prvo uključi tranzistor $T_P$ negativnim impulsom iza diferencijatora.
\end{columns}
\end{frame}
